\begin{workedexample}[Worked Example: Spin echo versus gradient echo]
A spin echo uses a $90^{\circ}$ excitation and a $180^{\circ}$ refocusing
pulse that cancels static field inhomogeneity, so its contrast follows true
$T_2$. A gradient echo omits the $180^{\circ}$ pulse and refocuses with gradients
alone --- faster (short TR) but sensitive to $T_2^{*}$. Long TR with short TE gives
proton-density weighting; short TR emphasizes $T_1$; long TE emphasizes $T_2$.
\end{workedexample}

\begin{keytakeaways}
\begin{itemize}[leftmargin=1.4em,itemsep=2pt]
\item Spin echo (with a $180^{\circ}$ pulse) images true $T_2$; gradient echo images $T_2^{*}$ and is faster.
\item Contrast: short TR $\to T_1$; long TE $\to T_2$; long TR/short TE $\to$ proton density.
\item Each excitation fills lines of k-space; the trajectory sets speed and artefacts.
\end{itemize}
\end{keytakeaways}

\begin{exercises}
\item Which sequence refocuses static field inhomogeneity: SE or GRE?
\item Which weighting comes from a short TR?
\item What does the $180^{\circ}$ refocusing pulse cancel?
\end{exercises}

\furtherreading{Nishimura~\cite{nishimura} (ch.~6); Haacke~\cite{haacke}.}
